% Chapter 12 — Presentation and Reporting Projections (MINSKY, 4 pp)
\section{Presentation and Reporting Projections}\label{ch:reporting}

This chapter discharges \S 1 and \S 3 of the constitution --- a report is a
projection of the one record, computed on demand --- and the auditability and
reproducibility commitments of \S 2.

\subsection{The projection principle}

Every reporting figure is a projection: a pure, total, deterministic function of
the record and of the declared terms the record carries, and of nothing else.
There is no reporting store and no parallel reporting pipeline; a report is
computed when asked for and discarded when read.

The principle is forced, not chosen. \S 1 removes reconciliation by recording
each fact exactly once; a reporting database would be a second store of the same
facts, maintained by its own logic on its own schedule, and two stores of one
fact eventually disagree. The divergence would surface in the documents third
parties read --- the balance sheet, the disclosure, the supervisory return ---
which is the worst possible place for the failure mode the architecture exists
to abolish. So no reporting figure is ever kept: what can be computed from the
coordinates, the units, and the declared terms is computed from them, every
time.

Three commitments of \S 2 are discharged at once, by construction rather than by
apparatus. \emph{Auditability}: a projection is a fold over the immutable log
(Chapter~2), so every figure decomposes mechanically into the transactions that
entered the fold --- drill-down is the fold's own trace. \emph{Reproducibility}:
the function is pure and its inputs are recorded, so any party holding the
record recomputes the same figure to the last minor unit; where a figure
involves value, the prices come from the pricing layer, itself reproducible from
recorded inputs (Chapter~16). \emph{Time travel}: a report as of any past date
is the same function applied to the log up to that date, and a restated report
is a new projection of a longer log containing recorded corrections --- never an
edited document.

One input class stands outside the record: reference data --- identifiers,
classifications, counterparty attributes --- owned by external authorities, per
the constitution's scope. It is supplied at the boundary by the authority that
owns it and joined to the projection as the report is built. The join attaches
labels and groupings; it can never alter a quantity. Gating admission and altering a
quantity are different acts, and neither contradicts the other: an eligibility predicate
may read declared or reference data to refuse a transaction at the door (Chapter~9), but a
refusal admits or rejects, it never rewrites a booked number --- and the report-time join
does not even gate, it only labels. A figure carrying a wrong
identifier is a labelling defect at the boundary, repaired by re-joining fresh
reference data to the unchanged projection --- never a ledger defect. The same
division of labour governs format: the ledger computes what the figures are; a
submission gateway, outside the boundary, arranges them into whatever shape a
supervisory template requires. Per-regime field inventories therefore live
outside this specification and are recorded as exclusions in the Exclusions
Register; where a template requires a figure, the figure is derived here from
the record and shown to satisfy the template --- the template is never the
derivation.

\begin{lstlisting}
-- One arrow, no state of its own, total in every argument.
-- Declared terms and conventions travel inside the prefix.
report :: LogPrefix -> BoundaryRef -> Report
\end{lstlisting}

\subsection{The presentation projection}

Chapter~9's regime key decides what a collateral inflow writes; the
presentation projection recovers from those writes the figures that a balance
sheet or a supervisory template asks of a firm that finances, pledges, and
re-uses assets. It reads coordinates, units, and declared terms alone. Four
figures cover the surface.

\emph{Transferred assets and associated liabilities.} Under title-transfer
financing the poster's asset re-books to the taker; the poster holds a
claim-for-equivalent unit and the taker a repurchase or redelivery obligation
unit, both per governing agreement (Chapter~9). The projection maps each claim
held under a financing agreement to a transferred asset at fair value --- the
claim priced by the pricing layer, parameterised by the taker's state
(Chapter~7) --- and maps the paired obligation to its associated liability. The
presentation is complete within the ledger's fair-value scope.

\emph{Encumbrance, poster side.} Under pledge the asset never leaves the
poster's owned coordinate; its encumbered mass is the posted-as-collateral ray,
per agreement. Under title-transfer financing the encumbered position is the
claim held under the financing agreement. Unencumbered mass is owned less
posted, and it is non-negative because the coverage invariant (Chapter~14)
bounds posting by ownership --- the report cannot show negative free assets
because the door refused the state that would produce one.

\emph{Re-use, receiver side.} Mass held on the received-as-collateral ray is
available for re-use. Actual re-use is read from re-pledge moves, each of which
carries a source-agreement reference beside its governing reference
(Chapter~9); the source is an observable fact of the custody instruction, so
nothing fictitious is stored, and received, re-used, and available project as
per-(unit, agreement) actuals.

\emph{Fungible received mass.} Title-transfer received cash and securities land
on owned, commingled with proprietary holdings of the same unit. The record
cannot say which dollar was re-invested, because no bank statement can say it
either: a per-dollar ``actual'' over commingled mass would be a stored fiction.
Figures over this mass --- how much of the received value was re-used, and where
--- are attributions, computed by the declared convention of the next
subsection.

\begin{center}
\begin{tabular}{@{}ll@{}}
\toprule
Reported figure & Read from --- and from nothing else \\
\midrule
Transferred asset / associated liability & claims and obligation units under financing agreements \\
Encumbrance under pledge & posted-as-collateral ray of owned units, per agreement \\
Encumbrance under financing & claims held under financing agreements \\
Re-use actuals (pledge form) & source-agreement references on re-pledge moves \\
Re-use of received fungibles & owned mass, obligation units, the declared convention \\
\bottomrule
\end{tabular}
\end{center}

\subsection{The attribution convention is total}

The convention is data, not code: it is declared once, in the declared-data
slot of Chapter~8, and consumed here --- behaviour carried as declared data,
per \S 2. It has exactly three components:

\begin{enumerate}[nosep]
\item the \emph{pool predicate} --- a declared wallet and move predicate
  defining the pool over which received value is attributed (for instance, the
  wallets and moves constituting the reinvestment book); it is total over recorded
  log state, ranging only over wallets and moves already on the log, so an input not
  on the record is not admissible and there is no hidden manual-tagging step outside
  the projection;
\item the \emph{period basis} --- the timestamp or period on which the pro-rata
  ratio is struck;
\item the \emph{rounding rule} --- with remainders booked explicitly, per
  \S 4's remainder discipline; nothing is silently dropped.
\end{enumerate}

\begin{principle}[Attribution totality]\label{prin:attr-total}
The declared convention is a total function of log state: for every log prefix
it assigns every attributable quantity; the attributed parts and the explicit
remainder partition the whole; and two readers of one log and one declared
convention compute identical figures to the last minor unit.
\end{principle}

A convention missing any of the three components is a defect: two readers of
one log would compute different figures, which is internal reconciliation
failure reintroduced through a report. The projection therefore fails closed
--- it refuses to produce the figure rather than improvise a ratio (\S 2,
Correctness).

\begin{lstlisting}
data Convention = Convention
  { pool     :: Predicate (Wallet, Move) -- who is in the pool
  , basis    :: PeriodBasis              -- when the ratio is struck
  , rounding :: RoundingRule }           -- remainder explicit
attribute :: Convention -> LogPrefix -> Partition Amount
\end{lstlisting}

\subsection{The reconciled surface}

An attributed figure is single-sided: it is the firm's own aggregate, and no
counterparty computes a matching number to break against. For such figures the
correctness standard is Principle~\ref{prin:attr-total} --- totality and
determinism --- not agreement with a second party, because no second party
exists.

One figure on this surface is genuinely two-sided: the collateralisation of net
exposure under a governing agreement. It is the specification's single reconciled
surface, present by design and named beside the projections, not a concession
extracted from them after the fact. The counterparty computes the same figure
from the same agreement, and the two are compared at the boundary. It is
therefore sourced from the governing agreement unit's declared terms ---
valuation percentages, thresholds, netting --- consistently with the
counterparty, and it is named beside the single-sided aggregates in every
presentation, so that a reader never mistakes an attribution for a reconciled
fact. Where the two sides disagree, the disagreement is a boundary
reconciliation item against an external authority --- the class the
constitution's scope leaves at the perimeter --- never an internal break: inside
the boundary there is one log and one declared convention.

\subsection{Episode: the pledged one-touch in the encumbrance projection}

OT-1, the one-touch of Chapter~2, pays \USD{1{,}000{,}000} on a touch of OMEGA
at or below 80.00; it is held by W-ALPHA, marked 400{,}000 pre-knock, and
pledged security-interest to counterparty B under agreement G at a declared
valuation percentage of 50\%: W-ALPHA posted $+1$, B posted $-1$ (Chapter~9).
Before the knock, the encumbrance projection reads, with no store of its own:
under G, mass posted 1; price 400{,}000; collateral value
$400{,}000 \times 50\% = \USD{200{,}000}$ against an exposure of
\USD{150{,}000} --- sufficiency holds.

\emph{The trigger.} OMEGA's official close 79.00 is recorded, at or below the
declared barrier 80.00; the Event Monitor emits the touch.

\emph{The contract.} OT-1's contract fires on the event with the current
projection of the ledger and proposes one moveless transaction.

\emph{The transaction.}
\begin{center}
\begin{tabular}{@{}rl@{}}
\toprule
1 & UnitStatus of OT-1: \emph{live} $\to$ \emph{triggered} \\
2 & conditional payment obligation created: writer owes W-ALPHA
    \USD{1{,}000{,}000}, predicate G's trapping terms \\
\bottomrule
\end{tabular}
\end{center}
No moves; no coordinate changes anywhere.

\emph{The door.} Conservation holds trivially --- there are no moves; writer
discipline --- the unit's contract is the sole writer of its UnitStatus; a
duplicate emission finds the unit already \emph{triggered} and proposes
nothing.

\emph{The read after.} The same projection, one transaction later: mass posted
under G is still 1 --- a lifecycle event extinguishes value, never mass
(\S 4) --- but the pricing layer, parameterised by unit state, now prices the
triggered OT-1 at 0, so collateral value reads $0 \times 50\% = \USD{0}$
against the unchanged exposure of \USD{150{,}000}. The shortfall does not live
in the report: it surfaces as G's margin-call obligation at G's declared
valuation watch, with a deadline (Chapter~9), and the report shows both the
collapsed line and the open obligation from the record alone.

\emph{The projection reads mass from the coordinates and value from the pricing
layer, so the collapse is on the report the instant it is on the record --- and
there was no report-side store to go stale.}

\subsection{Verification and the audit gate}

Per the Testability commitment, the projections of this chapter enter the
executable-check regime of Chapter~15 as properties over generated logs,
agreements, and conventions:

\begin{enumerate}[nosep]
\item \emph{Presentation completeness} --- every claim held under a financing
  agreement appears exactly once as a transferred asset, and its paired
  obligation unit exactly once as the associated liability.
\item \emph{Encumbrance bound} --- per (wallet, unit), reported encumbered mass
  never exceeds $\max(\text{owned}, 0)$: the coverage invariant of Chapter~14,
  read back from the report.
\item \emph{Attribution partition} --- attributed parts plus the explicit
  remainder equal the whole, for every generated log and declared convention.
\item \emph{Re-use bound} --- per source agreement, re-pledged mass never
  exceeds the mass received under it.
\item \emph{Mass--value separation} --- a lifecycle event inserted into a
  generated history changes reported values, never reported mass, exactly as
  the episode above exhibits.
\item \emph{Recomputation} --- every published figure equals the figure
  recomputed from the log prefix and the boundary join: the reporting analogue
  of \S 13's economic verification.
\end{enumerate}

Two readers of one log compute identical reports; the properties above make
that sentence executable rather than asserted. One gate remains before the
specification freezes: the auditor bench signs the faithfulness of the
presentation projection --- that the mapped figures state what a fair-value
presentation asks --- against these property tests. The sign-off is recorded
here as a Phase~4 gate of the drafting programme. It confirms that the
derivation is faithfully presented; the derivation itself rests on the
constitution alone.
